\begin{theorem}[Bridge from standard continuous characters]\label{mod:basic-standardcharacterbridge}
Let $F$ be a nonarchimedean local field.  If a standard continuous
additive character $\psi:F\to\mathbf C^\times$ is supplied together with
an integer $n$ and a proof that $n$ is its exact additive conductor, define
\texttt{continuousAddChar\_toData}$(F,\psi,n)$ by retaining $\psi$ as its
character field, $n$ as its conductor field, and the given exactness proof.
Thus evaluation of the wrapped character is literally evaluation of $\psi$,
and
\[
 \operatorname{AddTriv}(\psi,r)\quad\Longleftrightarrow\quad -n\le r.
\]
In particular, wrapping preserves the manuscript's additive sign convention
and does not weaken exact conductor to an upper bound.

Similarly, if a standard continuous quasi-character
$\chi:F^\times\to\mathbf C^\times$ is supplied with a natural number $m$
and a proof that $m$ is its exact multiplicative conductor, define
\texttt{continuousQuasiChar\_toData}$(F,\chi,m)$ without changing $\chi$.
Then
\[
 \operatorname{MulTriv}(\chi,r)\quad\Longleftrightarrow\quad m\le r.
\]
At the endpoint, conductor zero is equivalent to triviality on
$U_F^0=\mathcal O_F^\times$; it is not global triviality of $\chi$.

The compact unit group $U_F^0$ is compact as a subset of $F^\times$.
The continuous image of $U_F^0$ under $\chi$, and hence the norms of all
its values, is bounded.  Since every positive and negative power of a unit
remains in $U_F^0$, boundedness of these powers forces
\[
 \lVert\chi(u)\rVert=1\qquad(u\in U_F^0).
\]
This gives both the standard-character and packaged-data forms of compact
unit unitarity.

Finally, wrapping commutes with the manuscript's pullbacks.  If the trace
pullback $\psi\circ\operatorname{Tr}_{K/F}$, respectively norm pullback
$\chi\circ N_{K/F}$, is supplied with its exact conductor proof, its wrapped
character is exactly the trace, respectively norm, pullback of the wrapped
downstairs character.  Pointwise these identities are
\[
 \psi_K(x)=\psi(\operatorname{Tr}_{K/F}x),\qquad
 \chi_K(y)=\chi(N_{K/F}y),
\]
so the conversions needed by the final export do not alter either
mathematical character.
\LeanFile{LanglandsFirstMainLemma/Basic/StandardCharacterBridge.lean}
\lean{LanglandsFirstMainLemma.continuousAddChar_toData}
\lean{LanglandsFirstMainLemma.continuousAddChar_toData_character}
\lean{LanglandsFirstMainLemma.continuousAddChar_toData_conductor}
\lean{LanglandsFirstMainLemma.continuousAddChar_toData_apply}
\lean{LanglandsFirstMainLemma.continuousAddChar_toData_isConductor}
\lean{LanglandsFirstMainLemma.continuousAddChar_toData_trivialOnLattice_iff}
\lean{LanglandsFirstMainLemma.continuousQuasiChar_toData}
\lean{LanglandsFirstMainLemma.continuousQuasiChar_toData_character}
\lean{LanglandsFirstMainLemma.continuousQuasiChar_toData_conductor}
\lean{LanglandsFirstMainLemma.continuousQuasiChar_toData_apply}
\lean{LanglandsFirstMainLemma.continuousQuasiChar_toData_isConductor}
\lean{LanglandsFirstMainLemma.continuousQuasiChar_toData_trivialOnUnitFiltration_iff}
\lean{LanglandsFirstMainLemma.continuousQuasiChar_toData_conductor_eq_zero_iff}
\lean{LanglandsFirstMainLemma.continuousQuasiChar_norm_eq_one_of_mem_unitGroup}
\lean{LanglandsFirstMainLemma.LocalQuasiCharData.norm_character_eq_one_of_mem_unitGroup}
\lean{LanglandsFirstMainLemma.continuousAddChar_toData_compTrace_character}
\lean{LanglandsFirstMainLemma.continuousAddChar_toData_compTrace_apply}
\lean{LanglandsFirstMainLemma.continuousQuasiChar_toData_compNorm_character}
\lean{LanglandsFirstMainLemma.continuousQuasiChar_toData_compNorm_apply}
\leanok
\SourceText{Local notation, subsection Local fields, ideals, and characters; final export contract.}
\uses{mod:basic-characterconductors,mod:localfield-extension}
\FormalizationContract
\end{theorem}
